\newpage
\section{实训主要内容及过程}


\subsection{襄阳东发实习服务有限公司}
\subsubsection{实训预备}
必备物品：手记本、笔。

\subsubsection{实训内容}
\begin{enumerate}[label=\arabic*., labelindent=0pt, itemindent=2em, listparindent=0pt, left=1.5em]
\item \textbf{安全讲座 — 2025/09/14 14:30}\\
主题：\textbf{安全第一，预防为主——实习行为规范}

\paragraph{进厂规范}
\begin{enumerate}[label*=\arabic*.]
  \item 必须经过入厂教育后才能进厂。
  \item 实习生整队进出厂，按标识行走，不得单独行动。
  \item 禁止带食物进厂，保持环境卫生。
  \item 不得在车间门口、草坪等区域聚集或坐歇。
  \item 不做与实习无关之事，如打扑克、吸烟等。
  \item 严禁使用照相机、摄像机、手机拍照。
  \item 不得在生产现场扎堆围观，以免影响生产。
  \item 不得带走工厂文件或产品，避免损坏或受伤。
  \item 未经同意不得进入未安排的车间。
  \item 必须戴实习帽，衣着整齐，禁止高跟鞋、短裤等。
  \item 严禁靠近运行设备，走指定安全通道。
  \item 禁止进入天车吊装区域，注意避让叉车等车辆。
\end{enumerate}

\paragraph{住宿行为规范}
\begin{enumerate}[label*=\arabic*.]
  \item 妥善保管贵重物品，如丢失自负。
  \item 禁止在寝室吸烟、点火，注意用电安全。
  \item 外出前关闭电器，合理使用空调与热水，超额自费。
  \item 禁止私拉电线或改动电源，若造成损失须赔偿。
  \item 洗浴注意水温与热水卡使用，避免烫伤与卡片损坏。
  \item 每天按时返回公寓，外出须报备。
  \item 不得私自留宿外人。
  \item 遵守学习场所纪律，禁止吃零食或闲聊。
  \item 晾衣物避免影响通道。
  \item 雨天雨伞需妥善摆放，不得阻碍通道。
  \item 床上用品、安全帽等物品需按规定归还。
\end{enumerate}

\item \textbf{工艺讲座 — 2025/09/14 15:00}

机器人是由多个部分协同构成的复杂系统，主要包括以下几个方面：
\begin{enumerate}
  \item \textbf{执行机构}：完成实际操作任务的部分，如机械臂、移动平台等。
  \item \textbf{驱动装置}：为执行机构提供动力，常见方式包括：
    \begin{itemize}
      \item 电驱动：响应快、控制方便。
      \item 液压驱动：功率大、结构紧凑、刚度好，但存在液体泄漏的风险。
      \item 气动驱动：结构简单、维护方便，但控制精度相对较低。
    \end{itemize}
  \item \textbf{感知反馈系统}：通过传感器获取环境和自身状态信息，为控制提供数据支撑。
  \item \textbf{控制系统}：协调各部分的运行，包含：
    \begin{itemize}
      \item \textbf{硬件部分}：工业控制板卡，包括主控单元、信号处理电路等。
      \item \textbf{软件部分}：控制算法，实现运动规划与状态调节。
    \end{itemize}
\end{enumerate}
综上，机器人通过执行机构、驱动装置、感知反馈系统与控制系统的有机结合，形成一个完整的自动化智能系统。

\item \textbf{工业机器人操作培训}
在之后的工业机器人示教操作中，我们观察到了一款 Takafuji 的六轴工业关节型机器人。它主要由以下三大部分组成：
\begin{enumerate}
  \item \textbf{机器人本体}：包含六个可独立运动的关节轴，能够实现复杂的空间运动与操作任务。
  \item \textbf{控制器}：负责接收指令并调度各个关节的协调运动，保证机器人按照预定逻辑安全、稳定运行。
  \item \textbf{示教器}：提供人机交互界面，操作人员可通过手动方式控制机器人运动，记录轨迹并存储任务程序。\
  在实际运行阶段，机器人即可依照示教器预设路径自主完成动作。
\end{enumerate}
该培训的重点在于让学员掌握工业机器人示教操作的基本流程，包括手动控制、轨迹录制与程序运行，从而为后续更复杂的自动化操作打下基础。

\item \textbf{6S基础培训}
6S 培训是工厂和企业常用的现场管理方法，包含以下六个方面：
\begin{enumerate}
  \item \textbf{整理（Seiri）}：区分必需与非必需品，将不必要的物品清除出现场，保持环境整洁。
  \item \textbf{整顿（Seiton）}：对留下的物品科学布局，做到定点、定容、定量，方便查找与使用。
  \item \textbf{清扫（Seiso）}：保持设备与环境的清洁，防止脏污和异常堆积，及时发现隐患。
  \item \textbf{清洁（Seiketsu）}：将整理、整顿、清扫制度化，建立标准和规范，确保持续保持。
  \item \textbf{素养（Shitsuke）}：培养员工遵守制度和规范的习惯，形成自觉的良好行为方式。
  \item \textbf{安全（Safety）}：重视安全生产，设置警示标识和防护措施，杜绝事故发生。
\end{enumerate}
此外，6S 推行中强调“三要素”：人人参与、持续改善、形成文化。设置方法包括区域划分、标识上墙和责任到人；标识方法则包括地面线条、物品标签与颜色管理。安全措施方面需注意佩戴防护用品、设置安全警示、规范操作流程等。
\end{enumerate}


\subsection{东风日产风神襄阳汽车有限公司}
\subsubsection{入厂准备}
入厂前需佩戴安全帽、反光背心，并在闸口上交手机，确保人身与信息安全。

\subsubsection{参观内容}
在参观过程中，我们深入了解了整车总装线。总装大致分为：车身入线——底盘装配——内饰与仪表——动力系统与发动机装配——车门/玻璃/天窗装配——质量检测（EOL, End-Of-Line）。

参观重点包括：
\begin{itemize}
  \item \textbf{吸盘式机械臂装配天窗工位}：机械臂通过视觉系统抓取天窗组件并完成上装，吸盘可适配不同天窗表面。工位配备力控/扭矩检验，装配完成后通过螺栓检测与防水密封胶点检测进行质量确认。
  \item \textbf{AGV 物流系统}：在总装线承担座椅、车门总成、小件料盘等中短距离运输，实现“点对点”拉动补料。AGV 与地面磁带或导航系统协同工作，并与 WMS/MES 系统接口，实现智能化物流管理。
  \item \textbf{质量检测与工艺保障}：采用红外检测技术对玻璃进行裂纹检测；自动化设备完成表面清理与边缘均匀投胶，确保密封性。下线前进行扭矩检测、气密/水密性检测、功能测试（灯光、仪表、CAN 通讯自检），并采样上传 MES 建档以便追溯。
  \item \textbf{企业文化墙及设备展示}：文化墙区域展示了 ABB 集团的 5CARA 机械臂与乘用车电气模型，反映出现代汽车生产广泛使用总线通信与高速网络。部分工位还配备无磁条导航 AGV、视频监控系统与安全标识设施，保障生产透明性与安全性。
\end{itemize}

\section{中印南方印刷有限司}
